\chapter{准备环境与编译 SDK}
\setchapterheader{第\thechapter 章\quad 准备环境与编译 SDK}

\begin{intro}
下载 SDK，准备 Ubuntu 与 Conda 依赖，生成 C++ 例程和 Python 扩展；另列出官方 ARM Docker 环境的使用方法。
\end{intro}
\section{获取 SDK 与模型}
从\href{https://github.com/Noetix-Robotics/noetix_sdk_bumi}{Bumi SDK 仓库}获取程序：
\par\noindent\begin{minipage}{\linewidth}
\begin{command}
cd ~
git clone https://github.com/Noetix-Robotics/noetix_sdk_bumi.git
cd ~/noetix_sdk_bumi
\end{command}
\end{minipage}\par

项目内下载目录名为 \code{noetix_sdk_bumi-main} 时，后续路径相应替换。关节模型 \code{Bumi_URDF.zip} 在官方手册的 URDF 小节，训练框架位于\href{https://github.com/Noetix-Robotics/noetix_bumi_lab}{noetix\_bumi\_lab}。

\begingroup\small\renewcommand{\arraystretch}{1.15}
\begin{longtable}{@{}P{0.33\linewidth}P{0.59\linewidth}@{}}
\toprule
\textbf{目录或文件} & \textbf{用途} \\
\midrule\endhead
\code{build.sh} & SDK 构建脚本 \\
\code{CMakeLists.txt} & 依赖查找、Python 扩展和可执行目标 \\
\code{config/} & DDS 与策略配置 \\
\code{dds/} & 按 CPU 架构区分的 DDS 库 \\
\code{examples/cpp/} & C++ 例程 \\
\code{examples/py/} & Python 例程 \\
\code{include/} & 控制器与公共结构体头文件 \\
\code{lib/} & SDK 原生库 \\
\code{policy/} & 策略模型 \\
\code{pybind/} & Python 绑定源码 \\
\code{build/} & 编译生成的扩展模块与例程 \\
\bottomrule\end{longtable}\endgroup

\section{Ubuntu 依赖}
官方编译环境说明采用 Ubuntu 22.04。先安装工具链和开发包：
\par\noindent\begin{minipage}{\linewidth}
\begin{command}
sudo apt update
sudo apt install -y build-essential git cmake \
  libeigen3-dev libfmt-dev pybind11-dev python3-dev \
  python3-pip libyaml-cpp-dev libglib2.0-dev \
  libboost-all-dev libgoogle-glog-dev \
  libprotobuf-dev protobuf-compiler
\end{command}
\end{minipage}\par

本地 SDK 使用 C++17，CMake 最低版本声明为 3.16，并查找 ONNX Runtime、Eigen3、Threads 和 pybind11；策略目标还链接 yaml-cpp。

官方环境文档提供了 x86\_64 的 CMake 3.26.3 包、ARM 的 3.31.8 安装器，以及 ONNX Runtime、glog、spdlog 等附件。需要自行构建这些依赖时，从\href{https://noetixrobotics.feishu.cn/docx/PP2BdXkSUoQ2VOxctJKcEdHanAg}{SDK 编译环境部署流程}获取对应架构的资源。先运行 \code{uname -m} 确认 CPU 架构，再选安装包。
\section{ONNX Runtime 开发库}
CMake 需要找到 ONNX Runtime 的开发安装。使用官方环境说明中的源码附件构建时，在解压后的源码目录执行：
\par\noindent\begin{minipage}{\linewidth}
\begin{command}
./build.sh --config Release --build_shared_lib --parallel 2
cd build/Linux/Release
sudo make install
sudo ldconfig
\end{command}
\end{minipage}\par

这里的 \code{build.sh} 属于 ONNX Runtime 源码目录。SDK 自身也有同名脚本，执行前先用 \code{pwd} 确认所在目录。

如果开发库已安装在自定义目录，配置 SDK 时用 \code{onnxruntime_DIR} 指向包含 \code{onnxruntimeConfig.cmake} 的目录。\code{python -m pip install onnxruntime} 用于 Python 推理依赖，C++ 构建继续使用开发库。
\section{选择 bumi\_deploy 并构建}
按照第三章创建环境并安装 Python 包后，在 SDK 根目录配置：
\par\noindent\begin{minipage}{\linewidth}
\begin{command}
conda activate bumi_deploy
cd ~/noetix_sdk_bumi
cmake -S . -B build \
  -DPYBIND11_FINDPYTHON=ON \
  "-DPython_EXECUTABLE=$(python -c 'import sys; print(sys.executable)')" \
  "-Dpybind11_DIR=$(python -m pybind11 --cmakedir)"
cmake --build build --parallel 2
\end{command}
\end{minipage}\par

这组参数把 pybind11 和 Python 都指向当前环境。若 \code{build/} 已由其他解释器配置，先把它改名保存，例如 \code{mv build build-before-conda}，再重新配置。
\begin{expected}
\code{build/} 中生成 \code{example_high}、\code{example_low}、\code{example_policy}、\code{example_media}，以及 \code{highcontrol_py}、\code{lowcontrol_py}、\code{mediacontrol_py} 对应的扩展模块。
\end{expected}
\section{检查导入并启动例程}
先检查 Python 和模块位置：
\par\noindent\begin{minipage}{\linewidth}
\begin{command}
python -c "import sys; sys.path.insert(0, './build'); import highcontrol_py; print(sys.executable); print(highcontrol_py.__file__)"
\end{command}
\end{minipage}\par

完成网络配置并准备好吊架后，按开发任务选一个例程运行：
\par\noindent\begin{minipage}{\linewidth}
\begin{command}
python examples/py/test_high.py
\end{command}
\end{minipage}\par

低层控制使用 \code{examples/py/test_low.py}，策略示例使用 \code{examples/py/test_policy.py}。C++ 的相应程序位于 \code{build/}，例如：
\par\noindent\begin{minipage}{\linewidth}
\begin{command}
sudo ./build/example_high
\end{command}
\end{minipage}\par

\section{使用官方 ARM Docker 镜像}
官方提供的 \code{noetix_dds.tar.gz} 面向 aarch64。解压后包含镜像 \code{noetix_dds.tar} 和 Compose 配置。导入镜像并创建工作目录：
\par\noindent\begin{minipage}{\linewidth}
\begin{command}
tar -zxvf noetix_dds.tar.gz
docker load -i noetix_dds.tar
docker images
mkdir -p /home/noetix/work
\end{command}
\end{minipage}\par

Compose 中保留 host 网络和工作目录映射：
\par\noindent\begin{minipage}{\linewidth}
\begin{command}
services:
  noetix_dds:
    image: noetix_dds:arm64
    container_name: noetix_dds
    network_mode: host
    volumes:
      - /home/noetix/work:/home/noetix/work:rw
    stdin_open: true
    tty: true
\end{command}
\end{minipage}\par

\par\noindent\begin{minipage}{\linewidth}
\begin{command}
docker compose -f noetix_dds.yaml up -d
docker ps
docker exec -it noetix_dds /bin/bash
\end{command}
\end{minipage}\par

把 SDK 放到映射的工作目录中，在容器内继续配置和编译。容器中的路径、Python 和开发库按容器环境检查；在主机创建的 Conda 环境由主机终端使用。
